\section{模仿学习 vs 强化学习：实际对比}

本讲由 Ashish Kumar 客座讲授，聚焦于 RL 在机器人领域（特别是移动和操作）中的实际应用。讲者首先对比了"实际有效的"模仿学习和强化学习。

\begin{importantbox}{IL vs RL 的实践特征}
\begin{itemize}
    \item \textbf{模仿学习}：
    \begin{itemize}
        \item 数据高效，适合能获取高质量演示的场景
        \item 性能上限受限于演示者的能力
        \item 难以超越人类水平
    \end{itemize}
    \item \textbf{强化学习}：
    \begin{itemize}
        \item 可以发现超越人类的策略
        \item 通常需要仿真环境（sim-to-real transfer）
        \item 奖励设计是关键挑战
        \item 训练周期长但推理高效
    \end{itemize}
\end{itemize}
\end{importantbox}

\subsection{何时选择 RL}

\begin{knowledgebox}{RL 的适用条件}
RL 在以下场景中相比 IL 更有优势：
\begin{itemize}
    \item 有高质量的\textbf{仿真器}可用
    \item 任务目标可以用奖励函数量化
    \item 需要\textbf{超越}人类演示者的水平
    \item 演示数据难以获取（如危险操作）
    \item 需要对环境变化具有\textbf{鲁棒性}
\end{itemize}
\end{knowledgebox}

\subsection{本章小结}

IL 和 RL 各有适用场景，实际系统中常常结合使用（如 IL warm-start + RL fine-tuning）。

